\subsection*{Identificación e instrucciones}
\textbf{Código:} CUE-01, versión 1.0. \quad
\textbf{Dirigido a:} clientes (PI-05). \quad
\textbf{Duración estimada:} 5 minutos.

Instrumento preparado, todavía sin aplicar. Antes de responder se explicará la actividad y
se solicitará el consentimiento correspondiente. No escriba nombres, teléfonos, números de
equipo ni montos. Puede omitir preguntas o terminar cuando desee. Piense en su servicio más reciente.

\subsection*{Experiencia del servicio}
Marque una respuesta por fila: \textbf{Sí}, \textbf{En parte}, \textbf{No} o
\textbf{No aplica / no recuerda}. Las preguntas se refieren a lo que pudo conocer como cliente,
no a los registros internos del taller.

\medskip
\noindent\begin{tabularx}{\linewidth}{@{}L{0.9cm}Ycccc@{}}
\toprule
\textbf{ID} & \textbf{Pregunta} & \textbf{Sí} & \textbf{En parte} & \textbf{No} & \textbf{N/A} \\
\midrule
C1 & Al dejar el equipo, ¿pudo confirmar qué problema se registró? & $\square$ & $\square$ & $\square$ & $\square$ \\
C2 & Cuando consultó por el servicio, ¿recibió información sobre su estado? & $\square$ & $\square$ & $\square$ & $\square$ \\
C3 & Si se necesitó un trabajo adicional, ¿le pidieron autorización antes de hacerlo? & $\square$ & $\square$ & $\square$ & $\square$ \\
C4 & ¿Le avisaron cuando el equipo estuvo listo para retirar? & $\square$ & $\square$ & $\square$ & $\square$ \\
C5 & Al retirar el equipo, ¿quedó claro si el servicio estaba pagado o pendiente? & $\square$ & $\square$ & $\square$ & $\square$ \\
\bottomrule
\end{tabularx}

\subsection*{Comentario opcional}
\textbf{C6.} ¿Qué información le habría ayudado durante el servicio? No incluya datos personales.

\medskip\noindent\rule{\linewidth}{0.4pt}\par\medskip
\noindent\rule{\linewidth}{0.4pt}

\medskip
El equipo contará las respuestas por opción y separará las omitidas. No presentará este
cuestionario como una medición de tiempos de búsqueda o de pagos omitidos dentro del taller.
La trazabilidad de C1--C5 consta en la Tabla~\ref{tab:traza-instrumentos}; C6 ayuda a detectar
necesidades que deban consultarse antes de cambiar el alcance.
